\chapter*{Résumé}
\addcontentsline{toc}{chapter}{Résumé}

Le marché de l'importation en Algérie est aujourd'hui fragmenté, peu transparent et fortement basé sur des échanges informels, laissant place à des risques élevés de fraude. Ce projet de fin d'études présente \textbf{Sila DZ}, une plateforme intelligente d'importation qui met en relation directe les clients et les importateurs. 

Sila DZ digitalise et structure entièrement le processus d'importation à travers une application centralisée permettant la création d'offres, la gestion des commandes, la négociation et le suivi en temps réel. Grâce à un système de vérification des utilisateurs et à un modèle économique innovant basé sur des points et des abonnements, la plateforme vise à réduire les intermédiaires, faciliter l'accès à la micro-importation et renforcer la confiance. Ce mémoire détaille la conception, l'architecture logicielle, ainsi que l'implémentation de Sila DZ.

\vspace{1cm}

\textbf{Mots-clés :} Importation, Algérie, Plateforme B2B, Micro-importation, Confiance, Digitalisation.

\newpage
\chapter*{Abstract}
\addcontentsline{toc}{chapter}{Abstract}

The import market in Algeria is currently fragmented, opaque, and heavily reliant on informal exchanges, leading to high risks of fraud. This graduation project presents \textbf{Sila DZ}, a smart import platform that directly connects clients and importers. 

Sila DZ completely digitizes and structures the import process through a centralized application that allows offer creation, order management, negotiation, and real-time tracking. Through a user verification system and an innovative business model based on points and subscriptions, the platform aims to reduce intermediaries, facilitate access to micro-importation, and build trust. This thesis details the design, software architecture, and implementation of Sila DZ.

\vspace{1cm}

\textbf{Keywords :} Import, Algeria, B2B Platform, Micro-importation, Trust, Digitalization.
